% ID: SEL-D0067-OB-E01
% ESTADO: validada. S067.
\textbf{a)} Como $F=3(2x+y)-2$ y $2x+y\ge6$, se cumple $F\ge16$. Para alcanzar la cota, imponemos $y=6-2x$. Al sustituir en $2x+5y\le30$, resulta $x\ge0$. Al sustituir en $2x-y\le6$, resulta $x\le3$. Por tanto,
\[
\min F=\boxed{16},\qquad
\text{puntos de mínimo: }\boxed{\{(x,6-2x):0\le x\le3\}}.
\]
El conjunto de mínimos es todo el lado entre $(0,6)$ y $(3,0)$. La intersección de las otras dos fronteras es $(5,4)$, el tercer vértice del recinto.

\textbf{b)} Ambas respuestas señalan puntos de mínimo: $(1,4)$ está en el interior del segmento y $(3,0)$ es uno de sus extremos.
